\documentclass[12pt]{article}
\usepackage[utf8]{inputenc}
\usepackage[T1]{fontenc}
\usepackage[margin=1in]{geometry}
\usepackage{setspace}
\usepackage{xcolor}
\usepackage[colorlinks=true, urlcolor=blue!60!black, linkcolor=blue!60!black]{hyperref}
\setlength{\parindent}{0em}
\setlength{\parskip}{0.45em}
\renewcommand{\baselinestretch}{1.0}
\begin{document}
\begin{center}
{\Large \textit{Roe v. Wade}, 410 U.S. 113 (1973)\par}
\vspace{0.3em}
{\normalsize Case Brief --- One-Page Edition\par}
\end{center}

\textbf{Memory Jogger.} The abortion-privacy case and its trimester framework, modified in \textit{Casey} and overruled in \textit{Dobbs} (\href{https://supreme.justia.com/cases/federal/us/410/113/}{410 U.S. 113}).

\textbf{Facts.} Texas made abortion a felony unless necessary to save the mother's life. Norma McCorvey (``Jane Roe''), pregnant and unmarried in Dallas, sued District Attorney Henry Wade; she gave birth before the case ended.

\textbf{Procedural History.} A three-judge court declared the statutes void but denied an injunction; direct appeal.

\textbf{Votes.} 7--2. Blackmun wrote for the Court; Burger, Douglas, and Stewart concurred; White and Rehnquist dissented.

\textbf{Holding.} The claim was not moot. Due process privacy includes a qualified abortion right; a fetus is not a constitutional person; and the trimester framework governs, with post-viability bans allowed except to preserve the mother's life or health.

\textbf{Reasoning.} Restrictive bans were relatively recent and early abortion had become safe. Building on \textit{Griswold}, the Court located privacy in Fourteenth Amendment liberty and declined to decide when life begins. White called the ruling ``an exercise of raw judicial power''; Rehnquist invoked \textit{Lochner} and the laws of 1868.

\textbf{Examples.} \emph{Same-Side:} a first-trimester hospital-only rule. \emph{Opposite-Side:} barring non-physician abortions. \emph{Fence-Sitter:} a 24-hour waiting period.

\textbf{Critique.} Originalists and \textit{Dobbs} find no textual basis; Ely and Ginsburg faulted its reasoning.

\textbf{Key Quotation.} ``This right of privacy, whether it be founded in the Fourteenth Amendment's concept of personal liberty and restrictions upon state action, as we feel it is, or, as the District Court determined, in the Ninth Amendment's reservation of rights to the people, is broad enough to encompass a woman's decision whether or not to terminate her pregnancy.'' --- Justice Blackmun, Opinion of the Court, 410 U.S. at 153.

\end{document}
